% Generated draft: requires mathematical and source review.
\documentclass[11pt]{article}
\usepackage[margin=25mm]{geometry}
\usepackage{fontspec}
\setmainfont{Latin Modern Roman}
\usepackage{amsmath,amssymb,mathtools,unicode-math}
\setmathfont{Latin Modern Math}
\usepackage{enumitem,needspace,xurl,hyperref}
\setlength{\parindent}{0pt}
\setlength{\parskip}{5pt}
\setlength{\emergencystretch}{3em}
\newcommand{\N}{\mathbb{N}}
\newcommand{\Z}{\mathbb{Z}}
\newcommand{\Q}{\mathbb{Q}}
\newcommand{\R}{\mathbb{R}}
\newcommand{\C}{\mathbb{C}}
\newcommand{\F}{\mathbb{F}}
\newcommand{\A}{\mathbb{A}}
\newcommand{\PP}{\mathbb{P}}
\newcommand{\E}{\mathbb{E}}
\newcommand{\eps}{\varepsilon}
\newcommand{\dd}{\,\mathrm{d}}
\begin{document}
\section*{UnsolvedMath ID: 2658}
\textbf{Source problem code:} catalogue entry
\subsection{Definitions and mathematical statement}
For a group $G$ equipped with a specified order $\leq$, a group order automorphism is a bijection $\varphi:G\to G$ satisfying
\[
\varphi(xy)=\varphi(x)\varphi(y),\qquad
x\leq y\ \Longleftrightarrow\ \varphi(x)\leq\varphi(y).
\]
An inner group automorphism has the form $x\mapsto gxg^{-1}$ for one fixed $g\in G$. Under this group-order interpretation, the supplied existence question asks for a Dlab group with its intended order and an order automorphism $\varphi$ such that
\[
\forall g\in G\ \exists x\in G:\quad \varphi(x)\ne gxg^{-1}.
\]
Equivalently the order-preserving group automorphisms should contain an element outside $\operatorname{Inn}(G)$. No quotient by all inner automorphisms is presumed without knowing that the intended order makes them order preserving.

The record does not define a Dlab group or its distinguished order, nor specify whether that order is total or partial and what left/right invariance is required. It also does not explicitly say whether ``order automorphism'' must be a group automorphism preserving order or merely an automorphism of the underlying ordered set. The first displayed definition declares the usual ordered-group interpretation, rather than choosing an unsupported definition of the specialized group class. These omissions prevent a complete self-contained quantified formulation for that class; an authoritative Dlab-group definition is needed before selecting concrete examples or excluding them.

The source is Kourovka Notebook, New Problems, 21st issue (2026), Problem 21.149. Its background lists V. M. Kopytov and N. Ya. Medvedev and then A. V. Zenkov, without a theorem or complete bibliographic item. The supplied literature triage, checked 2026-08-17, labels the question open and reports neither a verified outer order automorphism nor a theorem excluding all of them. It records no distinct mathematical partial result beyond the ordered-group setting and suggests studying automorphism groups of concrete Dlab groups. The primary registry status is also ``open''. No such computation or new status assessment is made here.
\subsection{Short English statement}
Does some Dlab group have an order-preserving group automorphism that is not conjugation by a group element? The source records this as open but does not define the specialized group class or its order conventions, so those definitions must be recovered before a fully self-contained statement is possible.
\subsection{Sources}
\begin{itemize}
\item UnsolvedMath local registry, permanent ID 2658, source code catalogue entry, ``Kourovka Notebook Problem 21.149'', New Problems, 21st issue (2026). The background lists V. M. Kopytov, N. Ya. Medvedev and A. V. Zenkov without complete bibliographic details. Dataset attribution: Ulam AI and the UnsolvedMath Contributors; curation and metadata under CC BY 4.0.
\item Kourovka Notebook, 21st issue (2026), Problem 21.149, supplied tracker locator \url{https://www.kourovka-notebook.org/}. The local literature triage records checks on 2026-08-17.
\item No authoritative definition or primary reference for Dlab groups is supplied in this record; no missing book, article, theorem or URL is invented.
\end{itemize}
\end{document}
